\practicechapter{第 1 章实践：程序如何从源码变成进程}

本章对应第一册第 1 章“程序如何从源码变成正在运行的进程”。正文讲编译、链接、装载和进程，本章把这条路线压缩成一个可验证任务：编译一个最小 C++ 程序，再用 \code{cpu1ee\_cli} 和系统工具证明它确实变成了某个 ELF64 小端二进制。

\practicesection{一、对应原书问题：从源码到产物身份}

先写一个最小程序：

\begin{lstlisting}[language=C++]
int add_i32(int a, int b) {
    return a + b;
}

int main() {
    return add_i32(20, 22) == 42 ? 0 : 1;
}
\end{lstlisting}

这段代码本身不能被 CPU 直接执行。实践要补上的不是一句“编译器会处理”，而是完整证据：用什么命令编译，产物有多少字节，magic 是否为 ELF，machine 是否是 x86-64，entry address 是多少，哪些 section 出现，工具输出是否能复查同一字段。

\practicesection{二、从特例推到规律：magic 正确不等于解析安全}

最小反例是只有四个 magic 字节的文件：\code{0x7F 'E' 'L' 'F'}。它看起来像 ELF，但长度不足 ELF64 header。若解析器看到 magic 就继续读 machine 和 entry，就会越界读取。第二个反例是 ELF32 或大端 ELF。magic 仍然正确，但字段宽度和字节序不符合本卷第一版合同。

由此推出解析顺序：先读字节，计算 checksum；再检查 magic；再检查 header 长度；再检查 class 和 endianness；最后才读取 object type、machine 和 entry address。这个顺序在 \code{inspect\_elf} 里表现为一组提前返回。提前返回不是偷懒，而是在输入合同不满足时停止解释。

\practicesection{三、实践任务：编译、检查、交叉验证}

把上面的程序保存到临时目录后执行：

\begin{lstlisting}[language=bash]
g++ -std=c++20 -O0 -g hello.cpp -o hello-debug
g++ -std=c++20 -O2 hello.cpp -o hello-release

books/cpu-volume-1-practice/build/executable-evidence-debug/labs/executable_evidence/cpu1ee_cli \
  ./hello-debug
readelf -h ./hello-debug
\end{lstlisting}

记录 \code{cpu1ee\_cli} 输出中的 \code{bytes}、\code{checksum}、\code{machine}、\code{entry}、\code{sections} 和 \code{symbols}。再从 \cmd{readelf -h} 中找到 Class、Data、Machine、Entry point address。两个工具不需要输出格式一致，但字段含义必须能对应。

本章还要补一组坏输入测试思路。用空文件、四字节 magic 文件、把 fixture 的 class 改成 ELF32 的输入分别检查 diagnostic。源码里已经有 \code{RejectsNonElfMagic}、\code{RejectsShortElfHeader}、\code{RejectsUnsupportedClassOrEndian}，阅读它们并写出每个测试对应哪个失败特例。

\practicesection{四、验收：报告必须能复查}

本章交付物是 \filepath{reports/ch01-program-to-process.md}。通过标准如下：

\begin{enumerate}
  \item 报告写出编译命令，不只写“我编译了程序”。
  \item 报告保存 Debug 和 Release 两个产物的 byte size 与 checksum。
  \item 至少用 \cmd{readelf -h} 交叉验证 machine 和 entry 两个字段。
  \item 能解释为什么 entry address 不等于 \code{main}，并说明启动代码会先运行。
  \item 能把三个坏输入和三个 diagnostic 对上。
\end{enumerate}

完成本章后，读者不只是“知道程序会变成进程”，而是能拿出一个具体二进制，说明它从哪里来、是什么格式、当前解析器能证明什么。
